In this section, we show a bound on utility lost due to the agent having a larger discount factor than the correct one. Recall that $k = \lceil - \frac{1}{\lg \gamma}\rceil$. We use this here to simplify the bounds

In contrast with the other bounds, in this section, we do not use \Cref{lem:optimization_bound} and the proof directly gives a construction of a utility function and belief for which equality is achieved. 

\begin{repeatthm}{thm:imprecise_discount}
Let $\pi_{\gamma}$ and $\pi_{\gamma^*}$ be perfect expected utility maximizers with respect to discount rates $\gamma$ and $\gamma^*$ respectively for some $\gamma \leq \gamma^*$, either with or without the ability to self-modify.. Let $u$,$\rho$ be their utility function and belief. Then 
\[
\left|V^{\pi_{\gamma^*}}\left(\mae_{<t}\right)-V^{\pi_{\gamma}}\left(\mae_{<t}\right)\right| \leq \frac{{\gamma^*}^{k} + {\gamma^*}^{k-1} - 1}{1-\gamma^*} - {\gamma^*}^{k-1}\frac{\gamma^{k} + \gamma^{k-1} - 1}{\gamma^{k-1}(1-\gamma)} 
\]
\end{repeatthm}
\begin{proof}
The proof is the same for both agents with and without the ability to self-modify as it only compares the utilities gained by the respective agents and disregards what history lead to them.

The proof works by formulating the maximum possible amount of expected utility lost as a solution of an optimization problem which we then solve. The objective function of this optimization problem will be the amount of expected utility lost and the constraints will be chosen such that they exactly match the assumptions of the theorem.

Let $\hat{u_t}$ and $\hat{u}'_t$ be the expected utilities gained at timestep $t$ given the policies $\pi_{\gamma^*}$ and $\pi_{\gamma}$ respectively. We can now define
\begin{minipage}[t]{0.5\textwidth}
\begin{gather}
S_1 = \sum_{t=1}^\infty {\gamma^*}^{t-1} \hat{u}_t \\
S'_1 = \sum_{t=1}^\infty {\gamma^*}^{t-1} \hat{u}'_t \\
\end{gather}
\end{minipage}
\hspace{-3em}
\begin{minipage}[t]{0.5\textwidth}
\begin{gather}
S_2 = \sum_{t=1}^\infty \gamma^{t-1} \hat{u}_t \\
S'_2 = \sum_{t=1}^\infty \gamma^{t-1} \hat{u}'_t
\end{gather}
\end{minipage}

We know that the policy leading to utilities $\{\hat{u}'_i\}_i$ is optimal for $\gamma$. This can be written as $S_2 \leq S_2'$ or in other words
\[
\sum_{t=1}^\infty \gamma^{t-1} (\hat{u}_t - \hat{u}'_t) \leq 0
\]

We need to find the largest possible $\epsilon$ such that $S_1 = S_1' + \epsilon$ can be achieved, or equivalently, there exist $\{\hat{u}_i\}_i, \{\hat{u}'_i\}_i$ such that
\[
\sum_{t=1}^\infty {\gamma^*}^{t-1} (\hat{u}_t - \hat{u}'_t) = \epsilon
\]

Substituting $\Delta \hat{u}_t = \hat{u}_t - \hat{u}'_t$, we get the following optimization problem:
\begin{align}
\text{maximize }&\sum_{t=1}^\infty {\gamma^*}^{t-1} \Delta \hat{u}_t \\
\text{subject to }&\sum_{t=1}^\infty \gamma^{t-1} \Delta \hat{u}_t \leq 0, \\
&\Delta \hat{u}_t \in [-1, 1] \quad ,\forall t 
\end{align}

We now prove several lemmas about the structure of the solution to this optimization problem.
\begin{quotation}
\begin{lemma} \label{lem:structure}
The optimum is unique and there exists $k$ such that $\Delta \hat{u}_t = -1$ for $t < k$ and $\Delta \hat{u}_t = 1$ for $t > k$.
\end{lemma}
\begin{proof}
We show that if the second part holds, the optimum is unique. Suppose that there are two optima with $k_1$ and $k_2$ such that $k_1 < k_2$. However, this would mean that for every $i$, we would have that the first is elementwise less or equal to the second and for some element the inequality is strict. This would contradict optimality. For fixed $k$, the optimum choice of $\Delta \hat{u}_k$ is to set it as large as possible without violating constraints, meaning that the solution is unique. It remains to show existence of $k$ from the statement.

Let us fix $\Delta \hat{u}_t$ for all $t \not \in \{i,j\}$ for some $i < j$. We prove that if now $\Delta \hat{u}_i$ can be decreased and $\Delta \hat{u}_j$ increased without violating the second contraint, then the solution is not optimal. This means that in the optimum $\Delta \hat{u}_t$ have to be non-decreasing and at most one can be strictly between $-1$ and $1$, together implying the result.

We will be changing $\Delta \hat{u}_j$ and we find $\Delta_i$ such that the left-hand side of first constraint stays constant or, as the rest in fixed, $\gamma^{i-1} \Delta \hat{u}_i + \gamma^{j-1} \Delta \hat{u}_j$ stays constant. If we increase $\Delta \hat{u}_j$ by $\epsilon$, we have to decrease $\Delta \hat{u}_i$ by $-\gamma^{j-i} \epsilon$. The objective function then changes by $\epsilon {\gamma^*}^{j-1} - \epsilon \gamma^{j-i} {\gamma^*}^{i-1}$. We have
\begin{gather}
\epsilon {\gamma^*}^{j-1} - \epsilon \gamma^{j-i} {\gamma^*}^{i-1} = \epsilon {\gamma^*}^{i-1} ( {\gamma^*}^{j-i} - \gamma^{j-i} )
\end{gather}
Since we assume that $\gamma < \gamma^*$, this is positive, meaning that the objective function increased. Since all constraints remained satisfied, the solution was not optimal, a contradiction.
\end{proof}

\begin{lemma} \label{lem:tight_inequality}
The inequality (6) is tight in any optimal solution, that is $S_2 = S_2'$.
\end{lemma}
\begin{proof}
Suppose for contradiction that the inequality is sharp. This implies that there exists $i$, such that $\Delta \hat{u}_i < 0$. This means that we can increase the obective function by increasing $\Delta \hat{u}_i$ by an amount sufficiently small to not violate any constraint.
\end{proof}
\end{quotation}

Now that we know the structure of the optimal solution, we can get its value.

From \Cref{lem:tight_inequality}, we get
\[
\sum_{t=1}^\infty \gamma^{t-1} \Delta \hat{u}_t = 0
\]
Then we have
\begin{equation} \label{eq:what_is_delta}
- \gamma^{t-1} \Delta \hat{u}_k = \sum_{t=1}^{k-1} \gamma^{t-1} (-1) + \sum_{t=k+1}^{\infty} \gamma^{t-1} 1 = -\frac{1-\gamma^{t-1}}{1-\gamma} + \frac{\gamma^{t}}{1-\gamma} = \frac{\gamma^{t} + \gamma^{t-1} - 1}{1-\gamma}
\end{equation}

Using this, we can easily solve for $k$. By \cref{lem:structure} there is exactly one solution. The lemmas \ref{lem:structure} and \ref{lem:tight_inequality} then provide complete characterization of the optimal solution which allows us to compute the maximum possible $\epsilon$.

First of all, we find $k$. It is the first integer greater or equal to the solution $t$ to $\Delta \hat{u}_t = -1$.
\[
-1 = -\frac{\gamma^{t} + \gamma^{t-1} - 1}{\gamma^{t-1} (1-\gamma)}
\]
Solving this gives us $k = \lceil-\frac{1}{\lg \gamma}\rceil$.

Setting $\Delta \hat{u}_t = -1$ for $t < k$, $\Delta \hat{u}_t = 1$ for $t > k$ and using the above equality to get $\Delta \hat{u}_k$, we have
\[
\epsilon = \sum_{t=1}^{k-1} {\gamma^*}^{t-1} \Delta \hat{u}_t + \sum_{t=k+1}^{\infty} {\gamma^*}^{t-1} + {\gamma^*}^{k-1} \Delta \hat{u}_k = \frac{{\gamma^*}^{k} + {\gamma^*}^{k-1} - 1}{1-\gamma^*} - {\gamma^*}^{k-1}\frac{\gamma^{k} + \gamma^{k-1} - 1}{\gamma^{k-1}(1-\gamma)}
\]
\end{proof}
